Now let's apply information theory to time. We view a stochastic process as a chain of random variables over time:
\[ \dots, X_{-2}, X_{-1}, X_0, X_1, X_2, \dots \]
At any moment $t=0$, we stand between two worlds:
\begin{itemize}
    \item The \textbf{Past} (History): $\overleftarrow{X} = \dots X_{-2} X_{-1}$
    \item The \textbf{Future}: $\vec{X} = X_0 X_1 X_2 \dots$
\end{itemize}

\subsection{The Goal of Modeling}
Our goal is simple: predict the Future given the Past. We want to know the conditional distribution $P(\vec{X} | \overleftarrow{X})$.

However, the Past is infinitely long. We cannot just write down a lookup table for every possible infinite history. We need to \textit{compress} the past. We need to find a summary of the past that is just as good as the full past for predicting the future.

\subsection{Causal States}
Two different histories, $\overleftarrow{x}$ and $\overleftarrow{x}'$, might lead to the exact same predictions.
For example, in a periodic process "ABABA...", seeing "AB" implies the next symbol is "A". Seeing "ABAB" also implies the next symbol is "A". The histories are different, but their predictive implications are identical.

We say two histories are \textbf{causally equivalent} ($\sim_\epsilon$) if they have the same conditional distribution over futures:
\begin{equation}
    \overleftarrow{x} \sim_\epsilon \overleftarrow{x}' \iff P(\vec{X} | \overleftarrow{X} = \overleftarrow{x}) = P(\vec{X} | \overleftarrow{X} = \overleftarrow{x}')
\end{equation}

\begin{sidebar}{Occam's Razor}
This equivalence relation is the mathematical embodiment of Occam's Razor. It groups distinct histories into the \textit{same} state if they don't make a difference to the future. This ensures we never create a state unless it is theoretically necessary.
\end{sidebar}

\begin{intuition}
    If knowing you had coffee for breakfast changes your probability of being productive today, then "Coffee" and "No Coffee" are \textit{causally distinct} states. If the color of your socks does not change that probability, then "Red Socks" and "Blue Socks" are \textit{causally equivalent}.
\end{intuition}

The sets of equivalent histories are called \textbf{Causal States}, denoted $\mathcal{S}$.

\subsection{The Epsilon Machine}
The \textbf{$\epsilon$-machine} is the model constructed from these causal states. It consists of:
1. The set of Causal States $\mathcal{S}$.
2. The transition rules between them.

It is the unique, minimal, optimal predictor of the process.
